(1) 两支互感线圈的磁能为 \(U_L=\frac12L(i_1^2+i_2^2)+Mi_1i_2\)，电容能 \(U_C=(q_1^2+q_2^2)/(2C)\)。由 Lagrangian \(\mathcal L=U_L-U_C\) 得 \(L\ddot q_1+M\ddot q_2+q_1/C=0\)，\(M\ddot q_1+L\ddot q_2+q_2/C=0\)。磁能矩阵特征值 \(L+M,L-M\) 均为正需 \(|M|<L\)，若忽略此条件会出现负磁能和虚频率而非真实被动电感耦合。

(2) 取模态 \(q_\pm=(q_1\pm q_2)/\sqrt2\)，两运动方程相加或相减得到 \((L+M)\ddot q_++q_+/C=0\)、\((L-M)\ddot q_-+q_-/C=0\)。故 \(\omega_+=[C(L+M)]^{-1/2}\)，\(\omega_-=[C(L-M)]^{-1/2}\)。对 \(M>0\) 同向电流互感磁能更大，等效惯量为 \(L+M\)，频率较低；反向电流等效电感减小，频率较高。

(3) 给定初态投影为 \(q_\pm(0)=Q_0/\sqrt2\)，两模初速均零，所以 \(q_\pm(t)=Q_0\cos(\omega_\pm t)/\sqrt2\)。恢复原电荷 \(q_1=(Q_0/2)[\cos(\omega_+t)+\cos(\omega_-t)]\)，\(q_2=(Q_0/2)[\cos(\omega_+t)-\cos(\omega_-t)]\)。当频率相差很小时，第一回路振幅出现慢包络而第二回路出现互补包络；总能量 \(E=Q_0^2/(2C)\) 在两个模态和两电容—电感之间守恒，不能以 \(q_1^2+q_2^2\) 单项代替总能量。

(4) 令 \(\omega_0=1/\sqrt{LC}\) 且 \(|M|/L\ll1\)。二项展开给 \(\omega_+\simeq\omega_0[1-M/(2L)]\)，\(\omega_-\simeq\omega_0[1+M/(2L)]\)，故频率劈裂 \(\Delta\omega=\omega_--\omega_+\simeq\omega_0M/L\) 对 \(M>0\)；带符号 \(M<0\) 时按绝对值计拍频。在 \(M\to0\) 极限两模退化、\(q_1=Q_0\cos(\omega_0t)\)、\(q_2=0\)，说明两个回路真正解耦。
